\interlude{The Corpse of Milk}
One of the most prominent themes in Puella Magi Madoka Magica is decay. Entropy, obviously, is a form of decay, and thus the magical girls/witches are presented as a weapon against decay. However, there are other forms of decay at work: the city steadily degrades over the course of the series, from the bright clean spaces of Episode 1 to the crumbling ruins of Episode 12. Most notably, the mental states of the magical girls themselves decay. This is most pronounced with Sayaka's descent, but there's plenty of hints that the other magical girls suffer from severe depression, such as the fountain of what appears to be Prozac when Mami and Madoka have their heart-to-heart or the way Kyoko constantly eats her feelings. The entire point of the witch system is to get the magical girls to decay emotionally until they become witches; in a sense, all that Kyubey's system does is shift entropy from the physical decay of the universe into the emotional decay of the girls.

This constant presence of decay ties neatly into the series' Buddhist roots. The first of the Four Noble Truths (the core philosophical tenets of Buddhism) is the inevitability of dukkha, which translates roughly to suffering. There are three kinds of dukkha: the ordinary, obvious dukkha of illness, aging, and death; the anxious dukkha brought about by trying to hold on to things that are subject to time and therefore constantly changing; and the underlying dukkha inherent in all material things as a product of their transience \cite{ref108}. 

This last corresponds more-or-less directly to entropy, the principle that all material things must inevitably wind down \cite{ref109}. This inevitability of decay sounds like it ought to be a source of despair, but there are solutions. The primary Buddhist solution is detachment—to escape from this world is to escape the karmic cycle of inevitable despair \cite{ref110}. This is the door Madoka, in her role as the bodhisattva Kanon, opened for the magical girls at the end of the series. But is it the only solution? Is there no way to be happy within this transient world?

Western culture initially answers "no" as well. Christianity offers escape from this world to Heaven as its solution, with the added notion that at some future point God will destroy this world of suffering and replace it with a better one. However, in the Middle Ages and Renaissance a concept arose which gives an alternate path out of decay and despair: putrefaction.

Putrefaction is an alchemical concept, an alternate term for fermentation, but it came to refer to the way in which death and rot bring forth life \cite{ref111}. Consider a rotting piece of fruit. It is revolting to human senses, black and ugly and foul-smelling, but it is also a riotous explosion of new life such as mold and maggots. These in turn serve as nourishment for "higher" forms of life, until ultimately even the most exalted creatures depend on rot for their existence.

This is more than just the life cycle of biology, it is one of the most profound spiritual teachings of the alchemists: Death brings forth life. Rot and creation are one and the same. Decay is evolution.

Or put another way, flowers bloom in cemeteries. One such flower is the red spider lily, a crown of which adorns Homura's witch form in Rebellion. Frequently a vivid red color, unlike most flowers it loses its leaves before blossoming; it is associated with loved ones separated by fate and death and frequently planted in cemeteries in Japan \cite{ref112}. The connection to Homura's pain, separated from her beloved Madoka, is quite obvious.

However, the act of planting the flowers shows that one still acknowledges the lost loved one; love can endure where material existence has decayed away. Indeed, it is that love—originating in a destroyed universe—that brings Madoka back to Homura's illusory world. With her she brings two other beings, Charlotte and Sayaka. Both return out of duty and loyalty to Madoka, but later state additional motivations.

Unsurprisingly, given that she has been obsessed with cheese throughout Rebellion, Charlotte comes back for cheese. Cheese is an excellent symbol of putrefaction, being a delicious and nourishing substance that is at the same time essentially rotting milk. Charlotte is not alone in her motivations for return; all three of Madoka and her servants have returned for something valuable that emerged from decay. In the case of Madoka, it is her relationship with Homura, which evolved over the course of multiple timelines in which Madoka decayed from a bright, cheerful magical girl to the largely passive figure of the timeline showcased in the series, while Homura decays from timidity to being completely shut off. Sayaka, on the other hand, comes for her relationship with Kyoko, a relationship rooted in Kyoko's attempts to reach Sayaka when the latter's mental state was decaying rapidly.

The products of putrefaction, in other words, can be valuable. Decay is not an unmixed evil.

What, then, to make of the fact that Madoka's perfect, decay-free nirvana necessarily contains no cheese, both because there is no way of making it and because Charlotte would have no reason to leave otherwise? In a world free of decay, free of putrefaction, none of the beauty and life created by rot can exist. Neither cheese nor fire-forged friendships exist in Madoka's realm, so it cannot be considered an adequate solution to the problem of decay.

Only time and the likely inevitable sequels to Rebellion can tell if Homura's solution is any better.
